%% paper1.tex - SensorGuard-Research, Paper 1.
%% Springer LLNCS format (llncs.cls) for ICACS 2026 / Springer LNNS.
%% Compiles with pdfLaTeX on Overleaf. Bundle llncs.cls and splncs04.bst
%% with this file, and upload the figure PNGs from results/figures/ into the
%% project (or set \graphicspath below to their folder).
\documentclass[runningheads]{llncs}
\usepackage[T1]{fontenc}
\usepackage{cite}
\usepackage{amsmath,amssymb,amsfonts}
\usepackage{graphicx}
\usepackage{booktabs}
\usepackage{textcomp}
\usepackage{xcolor}
\usepackage{url}
\graphicspath{{figures/}{./}}

% Let floats settle near their references instead of piling up at the end:
% loosen LaTeX's default limits on how much of a page may be occupied by floats.
\renewcommand{\topfraction}{0.9}
\renewcommand{\bottomfraction}{0.8}
\renewcommand{\textfraction}{0.07}
\renewcommand{\floatpagefraction}{0.7}
\setcounter{topnumber}{3}
\setcounter{totalnumber}{4}

\begin{document}

\title{Comparative and Adaptive Evaluation of Unsupervised Machine Learning
Models for Multivariate Sensor Anomaly Detection in Industrial Systems}
\subtitle{A Cross-Dataset and Edge-Deployability Study}
\titlerunning{Adaptive Unsupervised Anomaly Detection for Industrial Sensors}

\author{Sarjan Shah \and Abhi Shah}
\authorrunning{S. Shah and A. Shah}
\institute{Department of Information Technology, Lok Jagruti University,
Vadodara, Gujarat, India \\
\email{shahsarjan968@gmail.com}, \email{abhi.shah@ljku.edu.in}}

\maketitle

\begin{abstract}
Industrial machinery increasingly carries many sensors whose readings can reveal
early faults, but labelled fault data is rarely available. This motivates
unsupervised anomaly detection, which learns only from normal operation. This
paper compares six unsupervised detectors (a statistical $z$-score baseline,
One-Class SVM, Local Outlier Factor, Principal Component Analysis (PCA)
reconstruction error, the Elliptic Envelope, and Isolation Forest) on two real
multivariate sensor benchmarks: the Skoltech Anomaly Benchmark (SKAB) and the
Secure Water Treatment (SWaT) testbed. The same leakage-free feature set and a
strictly chronological, validation-tuned protocol are used throughout. Alongside
F1, results are reported with the Matthews correlation coefficient (MCC) and
balanced accuracy, because F1 can reward a detector that simply flags almost
every reading. Three findings emerge. First, the best classical detector depends
on the dataset. Isolation Forest leads on SKAB, while the correlation-based
Elliptic Envelope leads on SWaT (MCC $0.77$), so a single winning model does not
transfer. A lightweight LSTM-autoencoder deep baseline, added for comparison,
sharpens this: it is the best detector on SWaT (MCC $0.84$) but falls below
chance on the smaller SKAB, so deep models do not remove the dataset dependence.
Second, a deployed static detector degrades under concept drift. On
SWaT, where normal operation drifts over time, a static model's false-alarm rate
on normal data rises from $8\%$ toward $100\%$. A lightweight, CPU-only drift
detector that retrains on recent data restores skill (MCC $0.24$ to $0.59$ on
SWaT), roughly halves false alarms, and keeps $90\%$ of attack recall. Third, the
full pipeline (the classical models and the drift adaptation, on both datasets)
runs
inside a container limited to Raspberry-Pi-class resources of one CPU core,
$1$\,GB RAM, and the ARM64 architecture. Inference stays below a millisecond per
reading and peak memory stays under $500$\,MB. The study contributes a
transparent multi-dataset baseline, evidence that drift adaptation is both needed
and feasible on commodity hardware, and a template for careful unsupervised
evaluation on real sensor data.

\keywords{anomaly detection \and unsupervised learning \and multivariate time
series \and concept drift \and edge computing \and predictive maintenance}
\end{abstract}

\section{Introduction}
Modern industrial systems such as pumps, motors, valves, and their connecting
pipework are routinely instrumented with sensors that measure vibration,
pressure, temperature, current, and flow. These signals carry early evidence of
developing faults. Detecting that evidence before failure is the goal of
condition monitoring and predictive maintenance, where machine-learning methods
have become central~\cite{rotating2023}. The main practical difficulty is that
labelled fault data is scarce. Faults are rare, varied, and expensive to induce
or annotate, so a supervised classifier trained on known faults tends to
generalise poorly to the next, unseen failure mode~\cite{chandola2009}.

Unsupervised anomaly detection avoids this problem. Instead of separating
``normal'' from labelled ``fault'' examples, an unsupervised (or novelty)
detector is trained only on data that represents healthy operation. At inference
time it flags readings that depart far enough from the learned notion of normal.
This setting matches industrial practice, where normal data is plentiful but
fault labels are not~\cite{belay2023}. Many such detectors exist, ranging from
simple statistical thresholds to deep reconstruction
networks~\cite{chalapathy2019}. For a practitioner the useful question is not
which method is most sophisticated, but which method actually performs best on
real sensor data under a fair evaluation, and whether it keeps working once
deployed.

That question is often answered on a single dataset, with optimistic evaluation
protocols, or without controlling for information leakage, all of which inflate
reported scores. It is also usually answered under the implicit assumption that
the data distribution stays fixed. Two consequences follow, and this paper
addresses both. A model declared best on one benchmark need not be best on
another. And a model that is deployed and then left unchanged can fail quietly as
the monitored system drifts away from its training distribution.

This paper makes the following contributions.
\begin{itemize}
\item It compares \textbf{six} classical unsupervised detectors, and a
lightweight LSTM-autoencoder deep baseline, on \textbf{two} real industrial
benchmarks (SKAB and SWaT) using one shared feature representation, so that the
comparison reflects the detector and not the preprocessing.
\item It follows a conservative protocol: chronological, per-experiment splits
with no shuffling; an unsupervised training set free of labelled anomalies; and
validation-based hyperparameter selection checked for ``flag-everything''
degeneracy. MCC and balanced accuracy are reported alongside F1 because they do
not reward that degeneracy, and a wall-clock feature leak found during the study
is documented and removed.
\item It shows that the \textbf{best detector depends on the dataset}, and that
this holds for deep models as well. Among the classical detectors, Isolation
Forest leads on SKAB and the correlation-based Elliptic Envelope leads on SWaT.
The deep baseline is the best detector of all on SWaT but falls below chance on
the smaller SKAB. A mechanistic explanation is given in terms of sensor count,
correlation, and data volume.
\item It reports a \textbf{concept-drift adaptation} result. A static detector
falls to chance-level skill under a regime shift on SKAB, and toward flagging all
normal data under genuine temporal drift on SWaT. A lightweight, CPU-only drift
detector that retrains on recent data recovers skill and roughly halves false
alarms.
\item It demonstrates \textbf{deployability}. The full pipeline runs within
Raspberry-Pi-class limits, measured by per-reading latency, model footprint, and
peak memory inside a constrained container.
\end{itemize}

The rest of the paper is organised as follows.
Section~\ref{sec:related} reviews related work. Section~\ref{sec:datasets}
describes the two datasets. Section~\ref{sec:method} covers the feature
engineering, models, evaluation metrics, drift adaptation, and constrained
runtime. Section~\ref{sec:results} presents the results.
Section~\ref{sec:discussion} discusses findings and limitations, and
Section~\ref{sec:conclusion} concludes.

\section{Related Work}
\label{sec:related}
Anomaly detection is a mature field with a well-established taxonomy of
statistical, distance-based, density-based, and classification-based
approaches~\cite{chandola2009}. Classical methods remain popular because they are
fast, easy to interpret, and need little tuning. Distance- and density-based
detectors such as the Local Outlier Factor (LOF) score a point by the relative
density of its neighbourhood~\cite{breunig2000}, extending earlier density-based
clustering work that separates dense regions from sparse, noise-like
points~\cite{ester1996}. Boundary methods such as the
One-Class SVM learn a surface that encloses the bulk of the normal
data~\cite{scholkopf2001}. Covariance-based methods model the normal data as one
distribution and flag points with large Mahalanobis distance; the Elliptic
Envelope uses the Minimum Covariance Determinant estimator so that a few
contaminating points do not distort the fit~\cite{rousseeuw1999}. Among
reconstruction-based detectors, PCA reconstruction error is the linear
archetype: it projects data onto a low-rank normal subspace and flags points
that reconstruct poorly~\cite{shyu2003}. Isolation Forest takes a different
route, isolating anomalies through random recursive partitioning on the premise
that anomalies are few and different and so are easier to isolate~\cite{liu2008}.

Deep learning has also been applied widely to anomaly
detection~\cite{chalapathy2019}. Sequence models such as LSTM encoder--decoders
reconstruct normal multivariate sequences and use the reconstruction error as an
anomaly score~\cite{malhotra2016}, extending earlier work that used plain
autoencoders as nonlinear reconstruction-based detectors~\cite{sakurada2014}.
These models capture temporal dependencies
that tabular methods miss, but they need much more data, compute, and tuning to
train and deploy. Recent benchmark studies report that, on real
multivariate-time-series data, deep detectors do not consistently outperform
well-tuned classical methods, and that careful feature engineering narrows or
removes the gap~\cite{belay2023,chalapathy2019}. This motivates the scope
adopted here. The study asks how far well-engineered features and standard
detectors can go on real benchmarks, and it includes a lightweight
LSTM-autoencoder as a deep baseline to measure that gap directly
(Section~\ref{sec:deepmethod}) rather than assuming it either way.

A second line of work concerns concept drift, meaning change in the data
distribution over time, and how to adapt to it through drift detection and
retraining on recent data~\cite{gama2014}, for example by tracking a statistic
over an adaptively resized window of recent observations~\cite{bifet2007}. Drift
matters for deployed
condition-monitoring systems, whose sensor baselines shift with wear and changing
operating conditions. This paper links the two threads. It first evaluates
classical detectors and then subjects the strongest models to drift, measuring
whether lightweight adaptation recovers performance.

\section{Datasets}
\label{sec:datasets}
Two public, real multivariate industrial sensor benchmarks are used.

\subsection{SKAB}
The Skoltech Anomaly Benchmark (SKAB)~\cite{skab2020} was recorded from a
water-circulation testbed built from a pump, pipework, and valves. Faults are
induced physically, mostly by partly closing a valve to restrict flow. Each
recording samples eight sensors at $1$\,Hz: two accelerometers (vibration RMS),
motor current, system pressure, water temperature, pump-housing thermocouple
temperature, supply voltage, and volumetric flow rate (RMS). The corpus has $35$
CSV files. One is an anomaly-free baseline run, and $34$ are fault runs in three
groups: \texttt{valve1} ($16$), \texttt{valve2} ($4$), and \texttt{other} ($14$
miscellaneous faults). The corpus holds $46{,}860$ rows in total. Each fault run
carries a per-timestamp \texttt{anomaly} label. Each CSV is an independent
experiment with its own clock, so the runs are not one continuous stream.
Table~\ref{tab:skabhead} shows a short excerpt.

\begin{table}[htbp]
\centering
\caption{Excerpt from a SKAB fault run (four consecutive rows). Four of the eight
sensors are shown for space; each row also carries the binary \texttt{anomaly}
label ($0$ = normal, $1$ = fault).}
\label{tab:skabhead}
\setlength{\tabcolsep}{4pt}
\footnotesize
\begin{tabular}{lccccc}
\toprule
datetime & Accel1RMS & Current & Pressure & Temp. & anomaly \\
\midrule
2020-03-09 10:34:33 & 0.0271 & 0.871 & 0.055 & 75.50 & 0 \\
2020-03-09 10:34:34 & 0.0270 & 1.301 & 0.055 & 75.54 & 0 \\
2020-03-09 10:34:35 & 0.0268 & 1.285 & 0.383 & 75.66 & 0 \\
2020-03-09 10:34:36 & 0.0268 & 1.365 & 0.055 & 75.58 & 0 \\
\bottomrule
\end{tabular}
\end{table}

\subsection{SWaT}
The Secure Water Treatment (SWaT) testbed~\cite{goh2016,mathur2016} is a
full-scale six-stage water-treatment plant instrumented with $51$ sensors and
actuators sampled at $1$\,Hz. It is one of the largest public industrial control
benchmarks, with about $1.44$ million rows in total: a long run of normal
operation followed by a run with $36$ labelled attacks that physically manipulate
the process. Of the $51$ channels, $25$ are continuous sensors (flows, tank
levels, pressures, chemical analysers) and $26$ are discrete actuator states
(valves, pumps). Table~\ref{tab:swathead} shows a short excerpt. The
fully-recorded normal block used here has $395{,}298$ rows. A further block, in
which seven channels were not recorded, is dropped rather than imputed. The
attack run contributes $54{,}621$ labelled rows. Compared with SKAB ($8$
sensors), the much larger and strongly inter-correlated SWaT sensor set turns out
to drive the cross-dataset findings in Section~\ref{sec:results}.

\begin{table}[htbp]
\centering
\caption{Excerpt from the SWaT run (four consecutive rows). Four of the $51$
channels are shown for space ($25$ continuous sensors are used in this study).
The sampling instant is given as a relative time index $t$ in seconds.}
\label{tab:swathead}
\setlength{\tabcolsep}{5pt}
\footnotesize
\begin{tabular}{cccccc}
\toprule
$t$ (s) & FIT101 & LIT101 & AIT201 & FIT401 & label \\
\midrule
0 & 2.494 & 817.67 & 262.63 & 1.712 & Attack \\
1 & 2.536 & 817.95 & 262.63 & 1.712 & Attack \\
2 & 2.536 & 817.95 & 262.63 & 1.712 & Attack \\
3 & 2.570 & 818.46 & 262.63 & 1.712 & Attack \\
\bottomrule
\end{tabular}
\end{table}

\section{Methodology}
\label{sec:method}

\subsection{Feature Engineering}
Raw sensor values discard temporal context, so each timestamp is augmented with
short-horizon causal statistics. For every continuous sensor, and using only the
current and past readings, the feature set adds a $5$-second rolling mean, a
$5$-second rolling standard deviation, a $30$-second rolling mean, and the first
difference. The raw value is kept, and one global feature,
\texttt{time\_since\_start}, is added. On SKAB this gives $41$ features
($8 \times 5 + 1$). On SWaT the same recipe over the $25$ continuous sensors
gives $126$ features. The discrete actuator channels are left out so that the
same rolling-statistics recipe applies to both datasets and the comparison stays
consistent.

Two choices protect the feature space. First, all rolling statistics are computed
within a single experiment segment. A window is never allowed to cross the
boundary between independent runs, or between a train and a test segment, which
would create a spurious link between unrelated conditions. Second, wall-clock
features such as hour-of-day and day-of-week are left out on purpose. Including
them gave a small score increase at first, but an ablation traced this to
leakage. Because the SKAB test set is the chronological tail of the same files
used for training, the hour within a short run is almost constant and acts as a
hidden file identifier. Dropping these features lowered the headline F1 but gave
a number that transfers to unseen data. All results use the leakage-free
representation.

\subsection{Splits and Unsupervised Training}
The data is a time series, so it is never shuffled.
\textbf{SKAB.} Each fault file is split chronologically into the first $80\%$ for
training and the final $20\%$ for testing, and these portions are concatenated
across files. The unsupervised training set combines the entire anomaly-free run
with the rows labelled normal (that is, with no fault present) from each training
portion, which keeps training free of labelled anomalies. This gives $28{,}368$ training rows and $7{,}503$ test
rows, of which $30.0\%$ are anomalous.
\textbf{SWaT.} The normal run is split chronologically into $70\%$ for training,
$15\%$ for validation-normal, and $15\%$ for test-normal, and the attack run is
split into validation and test halves. The model trains on normal rows only. This
gives $276{,}708$ training rows and an $86{,}606$-row test set, of which $31.5\%$
are anomalous.
On both datasets, hyperparameters are chosen on the validation split only, and
features are standardised with a scaler fit on the training set alone.

\subsection{Models}
Six unsupervised detectors are evaluated. All come from
scikit-learn~\cite{pedregosa2011} except the statistical baseline, and all sit
behind a common train/predict interface on the same scaled feature matrix, built
with NumPy~\cite{harris2020} and pandas~\cite{mckinney2010} for array and
tabular data handling. The
$z$-score baseline flags a row $\mathbf{x}$ when any standardised feature exceeds
a fixed threshold $\tau$,
\begin{equation}
\max_{j}\left| \frac{x_j - \mu_j}{\sigma_j} \right| > \tau,
\end{equation}
where $\mu_j$ and $\sigma_j$ are the training mean and standard deviation of
feature $j$. The One-Class SVM~\cite{scholkopf2001}
uses an RBF kernel. Local Outlier Factor~\cite{breunig2000} runs in novelty mode
so it can score held-out points. PCA reconstruction error~\cite{shyu2003} keeps
the components that explain $95\%$ of the training variance and scores by the
residual error. The Elliptic Envelope~\cite{rousseeuw1999} fits a Mahalanobis
model through the Minimum Covariance Determinant. Because the features are partly
collinear (rolling means track raw values) the raw covariance is singular, so the
data is PCA-whitened to $10$ components first, a standard way to keep the estimate
numerically stable. Isolation Forest~\cite{liu2008} scores by how shallowly a
point is isolated in an ensemble of random trees.

\subsection{Evaluation Metrics}
Precision, recall, and F1 on the anomaly class are reported together with the
confusion matrix. F1 is not used as the only headline metric. A detector that
alarms on almost every row reaches recall near $1$ and precision near the anomaly
base rate, which produces a middling F1 even though the detector separates
nothing. The Matthews correlation coefficient (MCC)~\cite{matthews1975} and
balanced accuracy are therefore reported and treated as primary, since both stay
informative under class imbalance where accuracy and F1 can be
misleading~\cite{powers2011}. From the
confusion-matrix counts,
\begin{equation}
\mathrm{MCC} = \frac{TP\cdot TN - FP\cdot FN}
{\sqrt{(TP{+}FP)(TP{+}FN)(TN{+}FP)(TN{+}FN)}},
\end{equation}
which is $0$ for a detector whose predictions are independent of the labels,
including one that flags every row, and $1$ for perfect separation. Balanced
accuracy is the mean of the per-class recall and equals $0.5$ in the same
flag-everything case. Both therefore reward only real separation of anomalies
from normal operation. This distinction becomes important in the concept-drift
analysis, where the static baseline turns into exactly that kind of
flag-everything predictor.

\subsection{Contamination Tuning}
The contamination parameter, meaning the assumed anomaly fraction, is tuned on
the validation split, and optima that correspond to flag-everything behaviour are
rejected. On SKAB, the Isolation Forest sweep has an interior optimum at $0.35$,
and this shared value is applied across the contamination/$\nu$-based models. On
SWaT, selection is by MCC. There, F1 keeps rising toward the degenerate
high-contamination end of the range, whereas MCC peaks at a selective value of
$0.05$ and then falls toward zero as the model degenerates, so MCC selection
lands on an operating point that can actually discriminate.

\subsection{Concept-Drift Detection and Adaptation}
\label{sec:driftmethod}
A streaming detector compares the feature distribution of an incoming batch
against a reference distribution using the two-sample Kolmogorov--Smirnov (KS)
$D$-statistic~\cite{massey1951}. For feature $j$, with empirical CDFs
$F^{\mathrm{ref}}_j$ and $F^{\mathrm{win}}_j$ over the reference set and the
current window, the per-feature statistic and the batch drift score are
\begin{equation}
D_j = \sup_{x}\left| F^{\mathrm{ref}}_j(x) - F^{\mathrm{win}}_j(x) \right|,
\qquad
\bar{D} = \frac{1}{d}\sum_{j=1}^{d} D_j,
\end{equation}
where $d$ is the number of features. Intuitively, the $D$-statistic measures how
far the current data has moved away from the training data on a scale from $0$ to
$1$, where $0$ means the two distributions are identical and $1$ means they do not
overlap at all. A rising score is therefore a simple, readable drift alarm. The
$D$-statistic is used rather than the KS $p$-value. With hundreds of rows per
feature in each batch, the $p$-value drops to almost zero for even a tiny change
and stops being informative, whereas the bounded $D$-statistic reports how large
the change is and does not grow just because there are more samples~\cite{gama2014}. When $\bar{D}$ passes a threshold, calibrated from ordinary
within-regime variation and never from the test stream, the model is retrained on
a sliding window of the most recent normal-labelled batches. Using a window rather than all past history
lets the model forget a superseded regime that would otherwise outnumber and
dilute the recent data.

Two drift scenarios are studied. On SKAB a regime shift is constructed. The
initial model is trained only on anomaly-free data plus \texttt{other}-family
normal data, and the test files are streamed in family order
(\texttt{other} to \texttt{valve1} to \texttt{valve2}), so \texttt{valve1} and
\texttt{valve2} arrive as unseen regimes. On SWaT the drift is real and temporal.
A detector trained on early normal data is streamed forward through later normal
operation, whose sensor baselines move over the multi-day run.

\subsection{Deep Sequential Baseline}
\label{sec:deepmethod}
To test whether the classical detectors' low compute costs a real accuracy
penalty, a lightweight LSTM~\cite{hochreiter1997} autoencoder is added as a deep
sequential baseline, implemented in PyTorch~\cite{paszke2019}. Unlike the tabular
detectors, it is given the temporal structure directly: its input is a sliding
window of $30$ consecutive raw (scaled) sensor readings ($8$ channels on SKAB,
$25$ on SWaT). A single-layer LSTM encoder with $64$ hidden units compresses the
window, a matching LSTM decoder reconstructs it, and the anomaly score for a row
is the mean squared reconstruction error of the causal window ending at that row.
Training uses normal windows only (Adam~\cite{kingma2015}, $25$ epochs), and the
decision threshold is selected on the validation split by MCC while rejecting
flag-everything cut-offs, the same rule used for the classical contamination. The
model is kept small (about $58$k parameters, $0.2$\,MB) so its footprint can be
compared against the classical detectors, and it is evaluated per row on the
identical SKAB and SWaT test splits.

\subsection{Constrained-Hardware Runtime}
To test edge deployability while using only commodity hardware, the pipeline runs
inside a Docker container limited to one CPU core, $1$\,GB of RAM, and the ARM64
architecture, which is the same instruction-set family as a Raspberry
Pi~\cite{merkel2014}. The measurements cover per-model training time, inference
latency in milliseconds per reading, on-disk model size, the on-device cost of a
drift check and a retrain, and whole-run peak memory. The effective CPU limit is
read from the container control group to confirm the single-core cap. This setup
reproduces a Pi's resource limits and instruction set, though not its slower
per-core speed, so the latency figures are an optimistic lower bound and are
reported as measured on a constrained ARM container rather than on a physical Pi.

\section{Results}
\label{sec:results}

\subsection{Model Comparison on SKAB}
Table~\ref{tab:skab} and Fig.~\ref{fig:skabcomp} give the comparison on SKAB.
Isolation Forest is the strongest detector on every metric (F1 $0.466$,
MCC $0.147$), with the Elliptic Envelope next. The absolute scores are modest.
The best precision is about $0.36$, which means roughly two of every three alarms
are false. That level is typical of unsupervised detection on real industrial
data, and it is reported as measured rather than inflated. The deep LSTM
autoencoder is the weakest entry here: its MCC of $-0.045$ is below chance and
below every classical detector, an outcome examined in
Section~\ref{sec:crossdataset}. Fig.~\ref{fig:skabconf} shows the confusion
matrices, and Figs.~\ref{fig:skabts}--\ref{fig:skabclose} illustrate Isolation
Forest's predictions against ground truth on one fault run.

\begin{table}[htbp]
\centering
\caption{Detector comparison on the SKAB test set, sorted by MCC (contamination
$0.35$ for the classical models). LSTM-AE is the deep sequential baseline. Best
per column in bold.}
\label{tab:skab}
\begin{tabular}{lccccc}
\toprule
Model & Precision & Recall & F1 & MCC & bAcc \\
\midrule
Isolation Forest   & \textbf{0.362} & \textbf{0.656} & \textbf{0.466} & \textbf{0.147} & \textbf{0.580} \\
Elliptic Envelope  & 0.342 & 0.539 & 0.418 & 0.086 & 0.547 \\
$z$-score          & 0.329 & 0.548 & 0.411 & 0.063 & 0.534 \\
One-Class SVM      & 0.326 & 0.537 & 0.406 & 0.055 & 0.530 \\
LOF                & 0.321 & 0.460 & 0.378 & 0.040 & 0.522 \\
PCA-recon          & 0.305 & 0.470 & 0.370 & 0.010 & 0.505 \\
LSTM-AE (deep)     & 0.251 & 0.127 & 0.169 & $-0.045$ & 0.482 \\
\bottomrule
\end{tabular}
\end{table}

\begin{figure}[htbp]
\centering
\includegraphics[width=\textwidth]{model_comparison.png}
\caption{Precision, recall, and F1 for the six models on the SKAB test set.}
\label{fig:skabcomp}
\end{figure}

\begin{figure}[htbp]
\centering
\includegraphics[width=\textwidth]{confusion_matrices.png}
\caption{Confusion matrices on the SKAB test set for all six models.}
\label{fig:skabconf}
\end{figure}

\begin{figure}[htbp]
\centering
\includegraphics[width=\textwidth]{timeseries_anomalies.png}
\caption{Isolation Forest predictions (orange) against ground-truth anomaly
regions (red shading) across the sensors of a SKAB fault run.}
\label{fig:skabts}
\end{figure}

\begin{figure}[htbp]
\centering
\includegraphics[width=\textwidth]{anomaly_closeup.png}
\caption{Close-up around the first changepoint of the run in
Fig.~\ref{fig:skabts}.}
\label{fig:skabclose}
\end{figure}

\subsection{Model Comparison on SWaT}
Table~\ref{tab:swat} and Fig.~\ref{fig:swatcomp} give the comparison on SWaT
(contamination $0.05$ for the classical models). The ranking changes in two ways.
Among the classical detectors, the correlation-based Elliptic Envelope now leads
by a wide margin (MCC $0.766$, F1 $0.839$), followed by the One-Class SVM, with
Isolation Forest, the SKAB winner, in third place. The deep LSTM autoencoder does
better still and tops the table (MCC $0.841$, precision $0.974$), the reverse of
its behaviour on SKAB. Absolute performance is much higher than on SKAB across
the board, which fits the richer sensor set providing stronger separation.

\begin{table}[htbp]
\centering
\caption{Detector comparison on the SWaT test set, sorted by MCC (contamination
$0.05$ for the classical models). The LSTM-AE deep baseline leads. Best per
column in bold.}
\label{tab:swat}
\begin{tabular}{lccccc}
\toprule
Model & Precision & Recall & F1 & MCC & bAcc \\
\midrule
LSTM-AE (deep)     & \textbf{0.974} & 0.805 & \textbf{0.881} & \textbf{0.841} & \textbf{0.897} \\
Elliptic Envelope  & 0.846 & 0.833 & 0.839 & 0.766 & 0.881 \\
One-Class SVM      & 0.653 & 0.896 & 0.755 & 0.634 & 0.838 \\
Isolation Forest   & 0.546 & 0.823 & 0.657 & 0.473 & 0.754 \\
$z$-score          & 0.349 & 0.955 & 0.511 & 0.180 & 0.567 \\
LOF                & 0.324 & \textbf{0.994} & 0.489 & 0.105 & 0.520 \\
PCA-recon          & 0.316 & 0.858 & 0.461 & 0.002 & 0.501 \\
\bottomrule
\end{tabular}
\end{table}

\begin{figure}[htbp]
\centering
\includegraphics[width=\textwidth]{swat_model_comparison.png}
\caption{Precision, recall, and F1 for the six models on the SWaT test set.}
\label{fig:swatcomp}
\end{figure}

\subsection{The Best Detector Is Dataset-Dependent}
\label{sec:crossdataset}
Fig.~\ref{fig:cross} places the two datasets side by side by MCC. The winner
changes. Isolation Forest is strongest on SKAB, and the Elliptic Envelope is
strongest on SWaT, there by a wide margin ($0.766$ against $0.473$), whereas the
two are close on SKAB. This is not accidental. It follows from the datasets. The
$51$ SWaT sensors move together in tight joint relationships, and an attack that
disturbs the process breaks those relationships even when no single sensor looks
extreme. That broken joint structure is exactly what a covariance-based detector
like the Elliptic Envelope reads directly. SKAB's $8$ sensors offer far less joint
structure to use, so the general-purpose partitioning of Isolation Forest is at
least as good. In practice a model called best on one benchmark should not be
assumed best on another, which makes multi-dataset evaluation necessary rather
than optional.

The deep LSTM autoencoder makes the same point more sharply. It is the best
detector of all on SWaT (MCC $0.841$), yet on SKAB it falls below chance
(MCC $-0.045$) and below every classical detector, even after training to
convergence on denser windows. The reason is again dimensionality and data
volume. On SWaT the model has $25$ correlated channels and hundreds of thousands
of normal windows, enough to learn a faithful temporal reconstruction so that
genuine attacks stand out as large errors: their median reconstruction error is
about three orders of magnitude above normal, cleanly separated by the threshold
(Fig.~\ref{fig:deeprecon}, right). On SKAB it has only $8$ channels and an order
of magnitude fewer windows, too little for the recurrent model to learn a normal
manifold that generalises, so the normal and anomaly reconstruction errors
overlap almost completely (Fig.~\ref{fig:deeprecon}, left). Training converged in
both cases (Fig.~\ref{fig:deeploss}), so the SKAB failure is not an artefact of
under-training; the loss plateaus and the error simply does not distinguish the
classes. The deep model does not remove the dataset dependence.
It amplifies it, and at a cost: it needs a deep-learning runtime far larger than
scikit-learn, benefits from a GPU for training, and its per-reading inference is
several times slower than the classical leaders (Section~\ref{sec:edge}), even
though its saved size ($0.2$\,MB) is comparable. Lightweight classical detectors
therefore remain the more robust and deployable choice across both datasets,
while a deep model pays off only where the data is high-dimensional and abundant
and the runtime budget allows it.

\begin{figure}[htbp]
\centering
\includegraphics[width=\textwidth]{cross_dataset_comparison.png}
\caption{Per-model MCC on SKAB versus SWaT, including the LSTM-autoencoder deep
baseline. The strongest model differs by dataset (Isolation Forest on SKAB, the
deep model on SWaT), and the deep model itself swings from best on SWaT to below
chance on SKAB.}
\label{fig:cross}
\end{figure}

\begin{figure}[htbp]
\centering
\includegraphics[width=\textwidth]{deep_recon_error.png}
\caption{Deep-model reconstruction error by class (log scale, box plots, fliers
hidden). On SKAB the normal and anomaly errors overlap, so the score cannot
separate the classes. On SWaT the anomaly errors sit far above the normal errors
and above the decision threshold, which is why the deep model leads there.}
\label{fig:deeprecon}
\end{figure}

\begin{figure}[htbp]
\centering
\includegraphics[width=\textwidth]{deep_training_loss.png}
\caption{LSTM-autoencoder training convergence (reconstruction MSE per epoch).
Both runs converge; the SKAB loss plateaus, confirming its poor detection is not
due to under-training.}
\label{fig:deeploss}
\end{figure}

\subsection{Concept-Drift Adaptation}
\label{sec:driftresults}
The comparisons above evaluate each detector once, on a fixed distribution. A
deployed detector meets data that drifts. This section subjects Isolation Forest
to drift and asks whether the lightweight adaptation of
Section~\ref{sec:driftmethod} counteracts the resulting degradation.

\paragraph{SKAB (regime shift)}
On the unseen valve regimes the static detector has no discriminative skill. It
reaches MCC $0.000$ and balanced accuracy $0.500$, no better than chance, and it
gets there by flagging almost everything, so its F1 of $0.43$ is an artefact of
the flag-everything behaviour that F1 fails to penalise. Retraining on drift
recovers real skill on those regimes (MCC $0.000$ to $0.106$, balanced accuracy
$0.500$ to $0.558$) and cuts the false-alarm load (precision $0.285$ to $0.336$).
Table~\ref{tab:drift} and Fig.~\ref{fig:skabdrift} summarise this.

\begin{table}[htbp]
\centering
\caption{SKAB drift: static vs.\ adaptive by regime (mean over batches). On the
unseen valve regimes the static model has zero skill (MCC $0.00$); adaptation
restores it.}
\label{tab:drift}
\begin{tabular}{llccc}
\toprule
Regime & Track & MCC & bAcc & F1 \\
\midrule
\texttt{other} (seen)   & static   & $-0.027$ & $0.484$ & $0.185$ \\
                        & adaptive & $-0.033$ & $0.485$ & $0.177$ \\
\midrule
valve (\textbf{unseen}) & static   & $0.000$ & $0.500$ & $0.430$ \\
                        & adaptive & $\mathbf{0.106}$ & $\mathbf{0.558}$ & $0.414$ \\
\midrule
whole stream            & static   & $-0.011$ & $0.493$ & $0.329$ \\
                        & adaptive & $\mathbf{0.049}$ & $\mathbf{0.528}$ & $0.317$ \\
\bottomrule
\end{tabular}
\end{table}

\paragraph{SWaT (genuine temporal drift)}
SWaT's normal operation drifts over its multi-day run. Fig.~\ref{fig:driftmotiv}
quantifies this. A detector trained on early normal data sees its false-alarm
rate on later normal data climb from $8\%$ toward $100\%$ as the sensor baselines
move, tracked by a rising KS drift score. Under the static-versus-adaptive stream
(Table~\ref{tab:swatdrift}, Fig.~\ref{fig:swatdrift}), retraining on drift more
than doubles skill (MCC $0.242$ to $0.594$; balanced accuracy $0.630$ to
$0.826$), cuts the false-alarm rate on normal data from $73\%$ to $24\%$, and
keeps most attack detection (recall $0.99$ to $0.90$). This drift is real rather
than constructed, and the gain is larger than on SKAB.

\begin{table}[htbp]
\centering
\caption{SWaT drift: static vs.\ adaptive (mean over $34$ streamed batches,
$25$ retrains). Adaptation halves false alarms and doubles skill while keeping
detection.}
\label{tab:swatdrift}
\begin{tabular}{lcc}
\toprule
Metric & Static & Adaptive \\
\midrule
MCC                         & $0.242$ & $\mathbf{0.594}$ \\
Balanced accuracy           & $0.630$ & $\mathbf{0.826}$ \\
Precision                   & $0.312$ & $\mathbf{0.583}$ \\
False-alarm rate (normal)   & $0.734$ & $\mathbf{0.244}$ \\
Recall (attacks)            & $0.993$ & $0.896$ \\
\bottomrule
\end{tabular}
\end{table}

\begin{figure}[htbp]
\centering
\includegraphics[width=\textwidth]{swat_drift_motivation.png}
\caption{SWaT normal operation drifts. The static detector's false-alarm rate on
normal data (red) climbs from $8\%$ toward $100\%$ as the KS drift score (blue)
rises, which motivates adaptation.}
\label{fig:driftmotiv}
\end{figure}

\begin{figure}[htbp]
\centering
\includegraphics[width=\textwidth]{swat_drift_adaptation.png}
\caption{SWaT concept-drift stream. Balanced accuracy over time for the static
(grey) and adaptive (blue) detectors, with dashed lines marking retrains. The
static detector decays toward the $0.5$ chance floor as normal drifts, while the
adaptive detector re-anchors after each retrain.}
\label{fig:swatdrift}
\end{figure}

\begin{figure}[htbp]
\centering
\includegraphics[width=\textwidth]{drift_adaptation.png}
\caption{SKAB concept-drift stream (balanced accuracy). Across the unseen
\texttt{valve1} and \texttt{valve2} regimes the static detector (grey) sits on
the $0.5$ chance floor, while the adaptive detector (blue) recovers skill.}
\label{fig:skabdrift}
\end{figure}

\subsection{Deployability on Constrained Hardware}
\label{sec:edge}
Table~\ref{tab:edge} and Fig.~\ref{fig:edge} report the pipeline run inside the
one-core, $1$\,GB, ARM64 container. The whole workload finishes with peak memory
at $484$\,MB, under half of the $1$\,GB cap, so it fits with room to spare.
Per-reading inference stays below a millisecond for every model and sits near two
microseconds for the two strongest ($0.0017$\,ms for the Elliptic Envelope,
$0.0019$\,ms for Isolation Forest). The leading models are also small on disk
($285$\,KB and $1.1$\,MB). The cost of adaptation is likewise low: a KS drift
check takes about $60$\,ms per batch and a full retrain about $0.11$\,s. A bounded
SWaT inference test inside the same container reproduces the cross-dataset
ordering (Elliptic Envelope MCC $0.81$ against Isolation Forest $0.66$), so the
constrained runtime does not change the result. These figures support edge
deployability, subject to the caveat in Section~\ref{sec:method} that a core
faster than a physical Pi makes the latency an optimistic bound. The deep LSTM
autoencoder is not included in the container run: its deep-learning runtime alone
exceeds the lightweight footprint measured here, and it benefits from a GPU for
training. On the host its saved model is small ($0.2$\,MB), but per-reading
inference is several times slower than the classical leaders. For constrained
deployment the classical detectors therefore remain preferable, the deep model
being justified only by the accuracy it buys on a high-dimensional dataset such
as SWaT.

\begin{table}[htbp]
\centering
\caption{Deployability inside a constrained ARM container (1 CPU, 1\,GB RAM):
per-reading inference latency and on-disk model size. Whole-run peak memory was
$484$\,MB of the $1024$\,MB cap.}
\label{tab:edge}
\begin{tabular}{lcc}
\toprule
Model & Inference (ms/reading) & Size (KB) \\
\midrule
$z$-score          & $<0.001$ & $1$ \\
Isolation Forest   & $0.0019$ & $1{,}134$ \\
Elliptic Envelope  & $0.0017$ & $285$ \\
PCA-recon          & $0.0132$ & $8$ \\
LOF                & $0.0901$ & $13{,}963$ \\
One-Class SVM      & $0.2415$ & $3{,}377$ \\
\bottomrule
\end{tabular}
\end{table}

\begin{figure}[htbp]
\centering
\includegraphics[width=\textwidth]{edge_benchmark.png}
\caption{Edge deployability inside a constrained ARM container (1 CPU, peak
$484$\,MB of $1024$\,MB). Per-model inference latency and on-disk size are shown
on a log scale. The leading detectors are both fast and compact.}
\label{fig:edge}
\end{figure}

\section{Discussion and Limitations}
\label{sec:discussion}
Three results stand out, and each is reported as measured, including where it
revises an earlier expectation.

The first is that no single model wins. Across the two datasets the best
classical detector changed, and the change tracks sensor count and correlation
structure rather than being incidental. The deep baseline behaves the same way,
only more extremely: best of all on SWaT, below chance on SKAB. This weakens any
universal best-model claim and argues for evaluation on more than one dataset.

The second is that static detectors fail quietly under drift. On both datasets a
static detector degrades, to chance on the unseen SKAB regimes and toward
flagging all normal data on SWaT, and F1 partly hides the failure. The
degradation is only fully visible under a metric that refuses to credit
flag-everything behaviour, which is itself a lesson about metric choice.
Lightweight, CPU-only adaptation counteracts the failure.

The third is that adaptation restores modest, not perfect, skill. On SKAB the
adaptive detector moves from no skill to some skill (MCC $0.00$ to $0.11$). On
SWaT the gain is larger (MCC $0.24$ to $0.59$) and false alarms roughly halve.
The value of the SKAB result is qualitative. The static model collapses to
exactly chance (MCC $0.00$), and a simple, label-free statistic, the KS
$D$-statistic, is enough to detect that collapse and trigger recovery
automatically. That mechanism is a structural foundation on which richer
retraining strategies can build, and the small absolute recovery on SKAB
reflects the difficulty of that particular regime-shift proxy rather than a limit
of the mechanism. The contribution is the demonstration that a static detector
fails under drift and that cheap adaptation counteracts it, not a claim of a
deployment-ready detector in absolute terms.

Several limitations bound these claims. Absolute performance on SKAB is modest,
with best precision near $0.36$, so the contribution is a fair comparison and an
adaptation mechanism, not a state-of-the-art detector. Both datasets are water
systems, so generalisation to other machine classes such as bearing wear,
turbomachinery, or electrical faults is not established. The SWaT data was
obtained through a public re-upload; the official iTrust source should be cited
and its terms observed for any formal release. The deep baseline tested here is a
single lightweight LSTM autoencoder~\cite{malhotra2016}; larger or better-tuned
sequence models might close the gap on SKAB or widen the lead on SWaT, and the
drift adaptation was studied only for the classical detectors, not the deep one.
The SKAB drift scenario is a constructed regime-shift proxy built from
fault families, while the SWaT drift is genuine; checking against progressive
physical wear is left for future work. Finally, the edge evaluation constrains
resources on an ARM host but is not a physical Raspberry Pi, so the latency is an
optimistic bound, and on-device timing on real hardware is future work.

\section{Conclusion}
\label{sec:conclusion}
This paper presented a controlled, reproducible, and deliberately broad study of
unsupervised anomaly detection for multivariate industrial sensors. Seven
detectors were evaluated: six classical models (a $z$-score baseline, One-Class
SVM, Local Outlier Factor, PCA reconstruction, the Elliptic Envelope, and
Isolation Forest) and one lightweight LSTM-autoencoder deep baseline. They were
tested on two real benchmarks of very different scale, SKAB (eight sensors, about
$47$k rows) and SWaT (fifty-one channels, about $1.44$ million rows), under a
single leakage-free feature set, a strictly chronological protocol, and five
metrics that include the discrimination-aware MCC and balanced accuracy. Every
detector was then examined under two further dimensions that a plain comparison
omits: distribution drift over time, and execution inside constrained edge
hardware.

Four concrete outcomes follow. First, the best classical detector depends on the
dataset. Isolation Forest leads on SKAB (MCC $0.147$) while the correlation-based
Elliptic Envelope leads on SWaT (MCC $0.766$), so a model declared best on one
benchmark does not transfer to another. Second, the deep baseline sharpens rather
than removes this dependence. It tops every detector on SWaT (MCC $0.841$) yet
falls below chance on the smaller SKAB (MCC $-0.045$), even after training to
convergence, which shows that a deep model repays its extra cost only where the
data is high-dimensional and abundant. Third, a deployed static detector degrades
under concept drift, reaching chance-level skill on unseen SKAB regimes and
driving its false-alarm rate on normal SWaT data from $8\%$ toward $100\%$. A
lightweight, CPU-only detector that senses drift with a Kolmogorov--Smirnov
statistic and retrains on recent data restores skill (SWaT MCC $0.242$ to
$0.594$), halves false alarms ($73\%$ to $24\%$), and keeps most attack recall
($0.99$ to $0.90$). Fourth, the whole pipeline, spanning the classical models and
the drift adaptation, runs inside a Raspberry-Pi-class container limited to one
CPU core, $1$\,GB of RAM, and the ARM64 architecture, with per-reading inference
below a millisecond and peak memory of $484$\,MB; the constrained run reproduces
the same cross-dataset ordering as the unconstrained host.

Taken together, these results contribute more than a single benchmark number.
They provide a transparent, multi-dataset and multi-model template for careful
unsupervised evaluation, and direct evidence that drift adaptation is both
necessary and feasible on commodity edge hardware. The immediate next step is
on-device validation on physical edge hardware such as a Raspberry Pi, to obtain
true bare-metal latency under real clock, thermal, and I/O constraints rather
than the optimistic bound reported here. Further work will extend the evaluation
to other machine classes such as rotating machinery and electrical faults, to
newer same-domain benchmarks such as WADI~\cite{ahmed2017}, apply the drift
adaptation to the deep model as well, and test it against progressive
physical wear.

\begin{credits}
\subsubsection{\discintname}
The authors have no competing interests to declare that are relevant to the
content of this article.
\end{credits}

\begin{thebibliography}{28}
\bibitem{rotating2023}
Das, O., Das, D.B., Birant, D.: Machine learning for fault analysis in rotating
machinery: a comprehensive review. Heliyon \textbf{9}(6), e17584 (2023)

\bibitem{chandola2009}
Chandola, V., Banerjee, A., Kumar, V.: Anomaly detection: a survey. ACM
Computing Surveys \textbf{41}(3), 1--58 (2009)

\bibitem{belay2023}
Belay, M.A., Blakseth, S.S., Rasheed, A., Salvo Rossi, P.: Unsupervised anomaly
detection for IoT-based multivariate time series: existing solutions,
performance analysis and future directions. Sensors \textbf{23}(5), 2844 (2023)

\bibitem{chalapathy2019}
Chalapathy, R., Chawla, S.: Deep learning for anomaly detection: a survey.
arXiv:1901.03407 (2019)

\bibitem{breunig2000}
Breunig, M.M., Kriegel, H.P., Ng, R.T., Sander, J.: LOF: identifying
density-based local outliers. In: Proceedings of the ACM SIGMOD International
Conference on Management of Data, pp. 93--104 (2000)

\bibitem{ester1996}
Ester, M., Kriegel, H.P., Sander, J., Xu, X.: A density-based algorithm for
discovering clusters in large spatial databases with noise. In: Proceedings of
the 2nd International Conference on Knowledge Discovery and Data Mining (KDD),
pp. 226--231 (1996)

\bibitem{scholkopf2001}
Sch\"olkopf, B., Platt, J.C., Shawe-Taylor, J., Smola, A.J., Williamson, R.C.:
Estimating the support of a high-dimensional distribution. Neural Computation
\textbf{13}(7), 1443--1471 (2001)

\bibitem{rousseeuw1999}
Rousseeuw, P.J., Van Driessen, K.: A fast algorithm for the minimum covariance
determinant estimator. Technometrics \textbf{41}(3), 212--223 (1999)

\bibitem{shyu2003}
Shyu, M.L., Chen, S.C., Sarinnapakorn, K., Chang, L.: A novel anomaly detection
scheme based on principal component classifier. In: Proceedings of the IEEE ICDM
Workshop on Foundations and New Directions of Data Mining, pp. 172--179 (2003)

\bibitem{liu2008}
Liu, F.T., Ting, K.M., Zhou, Z.H.: Isolation forest. In: Proceedings of the 8th
IEEE International Conference on Data Mining (ICDM), pp. 413--422 (2008)

\bibitem{malhotra2016}
Malhotra, P., Ramakrishnan, A., Anand, G., Vig, L., Agarwal, P., Shroff, G.:
LSTM-based encoder-decoder for multi-sensor anomaly detection. In: Proceedings
of the ICML Anomaly Detection Workshop (2016). arXiv:1607.00148

\bibitem{sakurada2014}
Sakurada, M., Yairi, T.: Anomaly detection using autoencoders with nonlinear
dimensionality reduction. In: Proceedings of the MLSDA 2014 2nd Workshop on
Machine Learning for Sensory Data Analysis, pp. 4--11 (2014)

\bibitem{gama2014}
Gama, J., \v{Z}liobait\.{e}, I., Bifet, A., Pechenizkiy, M., Bouchachia, A.: A
survey on concept drift adaptation. ACM Computing Surveys \textbf{46}(4), 1--37
(2014)

\bibitem{bifet2007}
Bifet, A., Gavald\`a, R.: Learning from time-changing data with adaptive
windowing. In: Proceedings of the 2007 SIAM International Conference on Data
Mining (SDM), pp. 443--448 (2007)

\bibitem{skab2020}
Katser, I.D., Kozitsin, V.O.: Skoltech Anomaly Benchmark (SKAB). Kaggle (2020).
\doi{10.34740/KAGGLE/DSV/1693952}

\bibitem{goh2016}
Goh, J., Adepu, S., Junejo, K.N., Mathur, A.: A dataset to support research in
the design of secure water treatment systems. In: Proceedings of the 11th
International Conference on Critical Information Infrastructures Security
(CRITIS), pp. 88--99 (2016)

\bibitem{mathur2016}
Mathur, A.P., Tippenhauer, N.O.: SWaT: a water treatment testbed for research
and training on ICS security. In: Proceedings of the International Workshop on
Cyber-physical Systems for Smart Water Networks (CySWater), pp. 31--36 (2016)

\bibitem{pedregosa2011}
Pedregosa, F., et al.: Scikit-learn: machine learning in Python. Journal of
Machine Learning Research \textbf{12}, 2825--2830 (2011)

\bibitem{harris2020}
Harris, C.R., et al.: Array programming with NumPy. Nature \textbf{585}(7825),
357--362 (2020)

\bibitem{mckinney2010}
McKinney, W.: Data structures for statistical computing in Python. In:
Proceedings of the 9th Python in Science Conference (SciPy), pp. 56--61 (2010)

\bibitem{matthews1975}
Matthews, B.W.: Comparison of the predicted and observed secondary structure of
T4 phage lysozyme. Biochimica et Biophysica Acta \textbf{405}(2), 442--451
(1975)

\bibitem{powers2011}
Powers, D.M.W.: Evaluation: from precision, recall and F-measure to ROC,
informedness, markedness and correlation. Journal of Machine Learning
Technologies \textbf{2}(1), 37--63 (2011)

\bibitem{massey1951}
Massey, F.J.: The Kolmogorov-Smirnov test for goodness of fit. Journal of the
American Statistical Association \textbf{46}(253), 68--78 (1951)

\bibitem{hochreiter1997}
Hochreiter, S., Schmidhuber, J.: Long short-term memory. Neural Computation
\textbf{9}(8), 1735--1780 (1997)

\bibitem{paszke2019}
Paszke, A., et al.: PyTorch: an imperative style, high-performance deep
learning library. In: Advances in Neural Information Processing Systems
(NeurIPS), vol. 32, pp. 8024--8035 (2019)

\bibitem{kingma2015}
Kingma, D.P., Ba, J.: Adam: a method for stochastic optimization. In:
Proceedings of the 3rd International Conference on Learning Representations
(ICLR) (2015)

\bibitem{merkel2014}
Merkel, D.: Docker: lightweight Linux containers for consistent development and
deployment. Linux Journal \textbf{2014}(239) (2014)

\bibitem{ahmed2017}
Ahmed, C.M., Palleti, V.R., Mathur, A.P.: WADI: a water distribution testbed for
research in the design of secure cyber physical systems. In: Proceedings of the
3rd International Workshop on Cyber-Physical Systems for Smart Water Networks
(CySWater), pp. 25--28 (2017)

\end{thebibliography}

\end{document}
